\documentclass{article}

\usepackage{amsmath,amssymb}
\usepackage{xurl}
\usepackage[hidelinks]{hyperref}

% Shared commands for notes in docs/paper-gaps/.

\DeclareUnicodeCharacter{2080}{\ifmmode{}_{0}\else\textsubscript{0}\fi}
\DeclareUnicodeCharacter{2081}{\ifmmode{}_{1}\else\textsubscript{1}\fi}
\DeclareUnicodeCharacter{2082}{\ifmmode{}_{2}\else\textsubscript{2}\fi}
\DeclareUnicodeCharacter{2099}{\ifmmode{}_{n}\else\textsubscript{n}\fi}

\newcommand{\Lean}{\textsc{Lean}}
\ExplSyntaxOn
\NewDocumentCommand{\leanid}{m}{
  \ifmmode
    \textsf{\detokenize{#1}}
  \else
    \group_begin:
    \urlstyle{sf}
    \tl_set:Nn \l_tmpa_tl {#1}
    \tl_replace_all:Nnn \l_tmpa_tl {\_} {_}
    \exp_args:NV \path \l_tmpa_tl
    \group_end:
  \fi
}
\ExplSyntaxOff
\providecommand{\tr}{\operatorname{tr}}
\newcommand{\tnleanIssueBase}{https://github.com/LionSR/TNLean/issues/}
\newcommand{\tnleanPrBase}{https://github.com/LionSR/TNLean/pull/}
\newcommand{\tnleanDocBase}{https://sirui-lu.com/TNLean/docs/}
\newcommand{\ghissue}[1]{\href{\tnleanIssueBase#1}{GitHub issue \##1}}
\newcommand{\ghpr}[1]{\href{\tnleanPrBase#1}{PR \##1}}
\newcommand{\doclink}[1]{\href{\tnleanDocBase#1.html}{\path{#1}}}

% Dirac notation and the identity operator, as in blueprint/src/macros/common.tex.
% \providecommand keeps note-local definitions authoritative.
\usepackage{dsfont}
\providecommand{\ket}[1]{|#1\rangle}
\providecommand{\bra}[1]{\langle #1|}
\providecommand{\braket}[2]{\langle #1 | #2 \rangle}
\providecommand{\ketbra}[2]{|#1\rangle\!\langle #2|}
\providecommand{\Id}{\mathds{1}}
\providecommand{\one}{\mathds{1}}
\providecommand{\C}{\mathbb{C}}
\providecommand{\MN}[1]{M_{#1}(\C)}
\providecommand{\MD}{M_D(\C)}

% Verdict marker: \gapnote{<kind>}{<status>}. Kind is the policy
% classification (clarification, local-correction, scope-restriction,
% unfaithful, false-source, open-gap); status is open, wip (elimination
% underway), resolved, or historical. Machine-read by the site index;
% prints a verdict line.
\newcommand{\gapnote}[2]{\par\noindent\textbf{Verdict:} #1 (#2).\par}


\title{The Printed One-Qubit $\mathbb{Z}_2\times\mathbb{Z}_2$ Symmetry:
the Anomaly Sign at $ab$}
\author{}
\date{2026-10-01}

\begin{document}
\maketitle
\gapnote{scope-restriction}{open}

\section{The source statement}

Ref.~\cite{GarreRubio2023MPOSym}, subsection ``MPO algebras representing
$\mathbb{Z}_2\times\mathbb{Z}_2$'', realizes one of the classes by the
operators $U_a=\prod_i CZ_{i,i+1}Z_i$, $U_b=\prod_i X_i$, and
$U_{ab}=U_aU_b$ on a periodic chain, and states: ``In this realization the
only non-trivial $3$-cocycle element is $\omega(ab,ab,ab)=-1$.''
In a normalized gauge the three diagonal values
$(\omega(a,a,a),\omega(b,b,b),\omega(ab,ab,ab))$ determine the class, so the
statement asserts the diagonal values $(+1,+1,-1)$.

\section{What is formalized}

For every choice of fusion tensors of the four printed tensors, the cyclic
product $\omega(ab,e,ab)\,\omega(ab,ab,ab)$ equals $-1$. Consequently the
class is nontrivial, and it has a normalized representative with
$\omega(ab,ab,ab)=-1$.\footnote{
    \leanid{KleinSymmetry.cyclicInvariant\_omega\_printedFamily},
    \leanid{KleinSymmetry.anomalyClass\_omega\_printedFamily\_ne\_zero}, and
    \leanid{KleinSymmetry.exists\_normalized\_omega\_printedFamily} in
    \doclink{TNLean/MPS/Examples/KleinSymmetryAnomaly}.}
The proof restricts the family to $\{e,ab\}$, where it is the decorated CZX
family.

\section{Missing piece and elimination plan}

The values at $a$ and at $b$ are not computed. The value at $b$ needs the
restriction to $\{e,b\}$, where all tensors have bond dimension one and the
cyclic product is one. The value at $a$ needs the anomaly of the order-two
family generated by the bond-two diagonal tensor of $U_a$. With both values,
the classification of Klein three-cocycles identifies the class as that of
$\omega_{001}=(-1)^{a_1b_2c_2}$.

\textbf{Verdict: scope restriction.} Only the sign at $ab$ of the printed
statement is formalized.

\bibliographystyle{alpha}
\bibliography{references}

\end{document}
